\clearpage
\begingroup
\fancyhead[L]{\small NACHLASS UND BRIEFWECHSEL}
\renewcommand{\headrulewidth}{0pt}
\par\medskip
\Needspace{8\baselineskip}
\pdfbookmark[2]{[5]}{pallas-IX-fragment5}
\begin{center}[5.]\end{center}
\begingroup
\small
\par\noindent\hspace*{\parindent}Für die Entfernung des Mars von der Pallas finden wir allgemein%
\HistoricalNoteEight{Die Zahlen in den Klammern \(\{\ \}\) sind Logarithmen.]}:
\[
\begin{aligned}
\rho\rho={}&rr+\{9{,}442\,6583-10\}r\sin(v+234^0\ 14'\ 37\overset{\prime\prime}{,}09)\\
&+2{,}321\,6384\\
&+\bigl(-0{,}432\,7895+\{0{,}473\,2062\}r\sin(v+54^0\ 14'\ 37\overset{\prime\prime}{,}09)\bigr)
  \mathord{.}\cos E'\\
&+\{0{,}407\,0481\}r\sin(v+332^0\ 32'\ 2\overset{\prime\prime}{,}60)
  \mathord{.}\sin E'\\
&+\{8{,}304\,6988-10\}\cos E^{\prime2}.
\end{aligned}
\]
\par\noindent\hspace*{\parindent}Dies gibt für \(M=30^0\) [--- für diesen Werth der mittlern Anomalie der
Pallas soll die Entwickelung von \(\frac{x'-x}{\rho^3}\) etc. nach Vielfachen von \(M'\) als
Beispiel gegeben werden; sie entspricht den mit \(M=30^0\) überschriebenen Verticalcolumnen
in den Tafeln auf den vorigen Seiten bis auf die dort hinzutretenden Grössen
\(\frac{x'}{r^{\prime3}}\) und \(\frac{y'}{r^{\prime3}}\)]:
\[
\begin{aligned}
{[1]}\quad\rho\rho={}&\{8{,}304\,6988-10\}\bigl(335{,}1002\\
&-\{2{,}502\,7828\}\cos(E'-198^0\ 39'\ 45\overset{\prime\prime}{,}90)+\cos E^{\prime2}\bigr)\\
{[}={}&\{8{,}304\,6988-10\}(A-B\cos(E'-C)+\cos E^{\prime2}).
\end{aligned}
\]
\par\noindent\hspace*{\parindent}Man bestimme nun die Wurzeln \(E'\) der Gleichung \(\rho\rho=0\), welche
\SourceLeaf{601}sämmtlich complex sind, und die dazu gehörigen Werthe von \(M'\) und bezeichne
mit \(\varphi'\) den reellen Theil desjenigen Werths von \(M'\), dessen imaginärer Theil der kleinste ist:]
\par\noindent\hspace*{\parindent}Zur Zerlegung der Grösse
\[
A-B\cos(E'-C)+\cos E^{\prime2}
\]
in die Factoren
\[
[2]\quad(M-\cos(E'-\varphi))\mathord{.}(N-\cos(E'+\varphi))
\]
hat man folgenden Gleichungen Genüge zu leisten:
\[
\begin{aligned}
{[3]}\quad MN&=A+\sin\varphi^2,\\
M\sin2\varphi&=B\sin(\varphi-C),\\
N\sin2\varphi&=B\sin(\varphi+C).
\end{aligned}
\]
In obigem Beispiele wird
\[
\begin{gathered}
\varphi=198^0\ 46'\ 43\overset{\prime\prime}{,}08,\\
\log M=0{,}023\,6709,\qquad
\log N=2{,}501\,6381.
\end{gathered}
\]
\par\noindent\hspace*{\parindent}Man setze nun ferner \(\frac1M=\sin\zeta\) [und man erhält für die gesuchte Wurzel]
\[
[4]\quad\cos(E'-\varphi)=\frac1{\sin\zeta},
\]
[woraus
\[
\sin(E'-\varphi)=i\operatorname{cotang}\zeta
\]
und]
\[
\begin{aligned}
E'-M'={}&[e'\sin E'=]e'\sin(E'-\varphi)\mathord{.}\cos\varphi\\
&+e'\cos(E'-\varphi)\mathord{.}\sin\varphi\\
={}&\frac{e'\sin\varphi}{\sin\zeta}
  +e'\cos\varphi\mathord{.}i\operatorname{cotang}\zeta.
\end{aligned}
\]
\par\noindent\hspace*{\parindent}[Also, indem man] \(\varphi-\frac{e'\sin\varphi}{\sin\zeta}=\varphi'\) [setzt,]
\[
[5]\quad M'-\varphi'=E'-\varphi-e'\cos\varphi\mathord{.}i\operatorname{cotang}\zeta.
\]
\par\noindent\hspace*{\parindent}[Die gesuchte Wurzel der Gleichung 4 hat offenbar
\begin{center}
\begin{tabular}{@{}ll@{}}
den reellen Theil: & \(\varphi\),\\
den imaginären Theil: & \(\operatorname{arc}\cos\frac1{\sin\zeta}\).\\
\end{tabular}
\end{center}
Der gesuchte reelle Theil von \(M'\) wird also nach 5 gleich \(\varphi'\).
\par\noindent\hspace*{\parindent}Aus 5 erhält man]
\par\noindent\hfill75\textsuperscript{*}\par
\SourceLeaf{602}
\[
\begin{aligned}
\cos(M'-\varphi')\bigl[={}&\frac1{\sin\zeta}
 \cos(ie'\cos\varphi\operatorname{cotang}\zeta)\\
&+i\operatorname{cotang}\zeta\mathord{.}
 \sin(ie'\cos\varphi\operatorname{cotang}\zeta)\bigr]\\
={}&\tfrac12c^{e'\cos\varphi\operatorname{cotang}\zeta}\operatorname{tang}\tfrac12\zeta\\
&+\tfrac12c^{-e'\cos\varphi\operatorname{cotang}\zeta}\operatorname{cotang}\tfrac12\zeta,
\end{aligned}
\]
[wo \(c\) die Basis der hyperbolischen Logarithmen ist.
\par\noindent\hspace*{\parindent}Setzt man]
\[
\operatorname{tang}\tfrac12\zeta\mathord{.}c^{e'\cos\varphi\operatorname{cotang}\zeta}
 =\operatorname{tang}\tfrac12\zeta',
\]
so ist
\[
[\cos(M'-\varphi')=\frac1{\sin\zeta'}
\]
und die Gleichung
\[
[6]\quad1-\sin\zeta'\cos(M'-\varphi')=0
\]
hat dieselbe in Frage kommende Wurzel, wie die Gleichung \(\rho\rho=0\). Man wähle also]
\(1-\sin\zeta'\cos(M'-\varphi')\) [als] Factor von \(\rho\rho\), nemlich
\[
\left(\frac1{\sin\zeta'}-\cos(M'-\varphi')\right)^{\frac32}=\lambda.
\]
\par\noindent\hspace*{\parindent}In unserm Beispiele:
\begin{center}
\begingroup
\footnotesize
\setlength{\tabcolsep}{1.5pt}
\renewcommand{\arraystretch}{1.18}
\begin{tabular}{@{}l@{\quad}l@{\quad}l@{}}
\begin{tabular}[t]{@{}l@{\ }r@{}}
\(M\)\enspace .\,.\,.\,.\,.\,.\,.\,.\,.\,. & \(0{,}023\,6709\)\\
\(e'\)\enspace .\,.\,.\,.\,.\,.\,.\,.\,.\,.\,. & \(4{,}283\,8772\)\\
\(\sin\varphi\)\enspace .\,.\,.\,.\,.\,.\,.\,. & \(9{,}507\,7380_n\)\\
\hline
\multicolumn{2}{@{}l@{}}{\(\log6535\overset{\prime\prime}{,}61=3{,}815\,2861_n\)}\\
\multicolumn{2}{@{}r@{}}{\(\varphi'=200^0\ 35'\ 38\overset{\prime\prime}{,}69\)}\\
\multicolumn{2}{@{}r@{}}{\(\zeta=71\ 15\ 15{,}57\)}\\
\end{tabular}
&
\begin{tabular}[t]{@{}l@{\ }r@{}}
\(\cos\varphi\)\enspace .\,.\,.\,.\,.\,.\,. & \(9{,}976\,2442_n\)\\
\(ke'\)\enspace .\,.\,.\,.\,.\,.\,.\,. & \(8{,}605\,0264\)\\
\(\operatorname{cotang}\zeta\)\enspace .\,.\,.\,. & \(9{,}530\,6736\)\\
\hline
& \(8{,}111\,9442_n\)\\
& \(-0{,}012\,9403\)\\
\(\operatorname{tang}\tfrac12\zeta\)\enspace .\,.\,.\,.\,. & \(9{,}855\,3053\)\\
\hline
\(\operatorname{tang}\tfrac12\zeta'\)\enspace .\,.\,.\,. & \(9{,}842\,3650\)\\
\multicolumn{2}{@{}r@{}}{\(\zeta'=69^0\ 38'\ 44\overset{\prime\prime}{,}22\)}\\
\end{tabular}
&
\begin{tabular}[t]{@{}l@{}}
\\
\\
\(\log\sin\zeta'=9{,}9719987\)\\
\(\displaystyle\frac1{\sin\zeta'}=1{,}0665993\)\\
\end{tabular}
\end{tabular}
\endgroup
\end{center}
\par\endgroup

\clearpage
\endgroup
